% 4.4 小群的小表示
\section{小群的小表示}

我们已经确定，按上一节的记号，$\mathbf{G}$ 的每个表示都要通过对 $\mathbf{K}^{i}$ 的某个表示 $\Gamma^{i}_{j}$ 施行诱导来得到。的确，$\mathbf{G}$ 的每个表示由两个标记刻画：$i$（星的标记）与 $j$（小表示的标记）。3.4 节描述的方法构成一个方案；现在来说明如何具体求出小表示 $\Gamma^{i}_{j}$。

由于 $\Gamma^{i}_{j} \downarrow \mathbf{T} = \lambda^{i}_{j}D^{i}$，表示 $\Gamma^{i}_{j}$ 限制到 $\mathbf{T}$ 上是 $\lambda^{i}_{j}$ 个 $D^{i}$ 的直和。因此选取合适的基后可以写
\begin{equation*}
\Gamma^{i}_{j}(t_{a}) = \begin{pmatrix} \mathbf{D}^{i}(t_{a}) & 0 & \dots\\ 0 & \mathbf{D}^{i}(t_{a}) & \dots\\ \vdots & \vdots \end{pmatrix},
\tag{4.4.1}
\end{equation*}
即 $\lambda^{i}_{j}$ 个矩阵 $\mathbf{D}^{i}(t_{a})$ 的直和。于是只需求出 $\gamma \in \overline{\mathbf{K}}^{i}$ 的矩阵代表 $\Gamma^{i}_{j}(r_{\gamma})$；因为此后 $\mathbf{K}^{i}$ 的所有元素的矩阵代表都可以由乘法规则构造出来。分两种情形讨论。

\subsection*{情形 (i)．$\mathbf{T}$ 为 Abel 群}

这涵盖空间群的情形，事实上也涵盖一切具有真不变子群、且该子群是循环群直积的有限群。此时 $D^{i}$ 是一维的，$\Gamma^{i}_{j}$ 的维数等于 $\lambda^{i}_{j}$。而且 $\Gamma^{i}_{j}(t_{a})$ 等于 $\lambda^{i}_{j}D^{i}(t_{a})\mathbf{1}$。

现在回忆对一切 $\beta, \gamma, \delta \in \overline{\mathbf{K}}^{i}$ 必须成立的下列关系：
\begin{equation*}
r_{\gamma}r_{\delta} = r_{\gamma\delta}t_{\gamma\delta},
\tag{4.4.2}
\end{equation*}
\begin{equation*}
t_{\beta\gamma, \delta}[t_{\beta\gamma}]_{\delta} = t_{\beta, \gamma\delta}t_{\gamma\delta}
\tag{4.4.3}
\end{equation*}
（见式 (4.2.3) 与 (4.2.11)）。于是在表示 $\Gamma^{i}_{j}$ 中必须有
\begin{equation*}
\Gamma^{i}_{j}(r_{\gamma})\Gamma^{i}_{j}(r_{\delta}) = \Gamma^{i}_{j}(r_{\gamma\delta})\Gamma^{i}_{j}(t_{\gamma\delta})
\tag{4.4.4}
\end{equation*}
\begin{equation*}
= \mathbf{D}^{i}(t_{\gamma\delta})\Gamma^{i}_{j}(r_{\gamma\delta})
\tag{4.4.5}
\end{equation*}
以及
\begin{equation*}
\mathbf{D}^{i}(t_{\beta\gamma, \delta})\mathbf{D}^{i}([t_{\beta\gamma}]_{\delta}) = \mathbf{D}^{i}(t_{\beta, \gamma\delta})\mathbf{D}^{i}(t_{\gamma\delta}).
\tag{4.4.6}
\end{equation*}
但 $D^{i} \simeq D^{i}_{\delta}$，因为 $D^{i}$ 是一维的且 $\delta \in \overline{\mathbf{K}}^{i}$（等价的一维矩阵相等）。于是式 (4.4.6) 成为
\begin{equation*}
\mathbf{D}^{i}(t_{\beta\gamma, \delta})\mathbf{D}^{i}(t_{\beta\gamma}) = \mathbf{D}^{i}(t_{\beta, \gamma\delta})\mathbf{D}^{i}(t_{\gamma\delta}).
\tag{4.4.7}
\end{equation*}
现在考虑群 $\overline{\mathbf{K}}^{i}$。令 $\Gamma^{i}_{j}(r_{\gamma}) = \Lambda^{i}_{j}(\gamma)$、$\mathbf{D}^{i}(t_{\beta\gamma}) = \mu^{i}(\beta, \gamma)$，则式 (4.4.5) 与 (4.4.7) 成为
\begin{equation*}
\Lambda^{i}_{j}(\gamma)\Lambda^{i}_{j}(\delta) = \mu^{i}(\gamma, \delta)\Lambda^{i}_{j}(\gamma\delta)
\tag{4.4.8}
\end{equation*}
与
\begin{equation*}
\mu^{i}(\beta\gamma, \delta)\mu^{i}(\beta, \gamma) = \mu^{i}(\beta, \gamma\delta)\mu^{i}(\gamma, \delta).
\tag{4.4.9}
\end{equation*}
式 (4.4.8) 与 (4.4.9) 表明，复数 $\mu^{i}$ 构成 $\overline{\mathbf{K}}^{i}$ 的一个射影表示的因子系统。而且该射影表示的维数为 $\lambda^{i}_{j}$ 且不可约，因为 $\Gamma^{i}_{j}(r_{\gamma}) = \Lambda^{i}_{j}(\gamma)$ 已不可约。

顺带地，我们验证了：由于 $\sum_{j}\left(\lambda^{i}_{j}\right)^{2} = |\overline{\mathbf{K}}^{i}|$，$\overline{\mathbf{K}}^{i}$ 的这些射影表示的维数平方和等于 $\overline{\mathbf{K}}^{i}$ 的阶。

当元素 $r_{\gamma}$ 可以选成构成一个群时会出现简化。此时对一切 $\gamma, \delta \in \overline{\mathbf{K}}^{i}$ 有 $t_{\gamma\delta} = t_{1} = e$。因子系统由全为 1 的数组成，射影表示成为普通矢量表示。对空间群而言这总会发生在 $\mathbf{k}$ 为布里渊区内点或空间群为点式的情形。

\subsection*{情形 (ii)．$\mathbf{T}$ 为非 Abel 群}

若 $D^{i}$ 是一维表示，则理论恰如情形 (i)。的确，情形 (i) 中用到 $\mathbf{T}$ 为 Abel 的唯一地方，是开头指出的那时 $D^{i}$ 必为一维。若 $D^{i}$ 的维数 $d_{i} > 1$，则下列事实成立：

（i）$\Gamma^{i}_{j}$ 的维数为 $d_{i}\lambda^{i}_{j}$，其中 $\Gamma^{i}_{j} \downarrow \mathbf{T} = \lambda^{i}_{j}D^{i}$。

（ii）代表 $\Gamma^{i}_{j}(t_{a})$ 的矩阵形式 (4.4.1) 仍成立，但 $\mathbf{D}^{i}(t_{a})$ 现在是 $d_{i} \times d_{i}$ 矩阵。

（iii）式 (4.4.5) 与 (4.4.7) 被替换为
\begin{equation*}
\Gamma^{i}_{j}(r_{\gamma})\Gamma^{i}_{j}(r_{\delta}) = \Gamma^{i}_{j}(r_{\gamma\delta})\Gamma^{i}_{j}(t_{\gamma\delta})
\end{equation*}
与
\begin{equation*}
\mathbf{D}^{i}(t_{\beta\gamma, \delta})\mathbf{D}^{i}_{\delta}(t_{\beta\gamma}) = \mathbf{D}^{i}(t_{\beta, \gamma\delta})\mathbf{D}^{i}(t_{\gamma\delta}),
\tag{4.4.10}
\end{equation*}
其中 $\Gamma^{i}_{j}(t_{\gamma\delta})$ 由式 (4.4.1) 给出，且 $\mathbf{D}^{i}_{\delta}(t_{\beta\gamma}) = \mathbf{D}^{i}(r^{-1}_{\delta}t_{\beta\gamma}r_{\delta})$。注意尽管 $D^{i}_{\delta} \equiv D^{i}$，这一等价一般不能取成恒等，因此 (4.4.10) 不再化为射影表示的乘法规则。

问题从这里开始不再容易。若 $\mathbf{G}$ 是半直积群，则对一切 $\gamma, \delta$ 有 $t_{\gamma\delta} = t_{1}$，式 (4.4.4) 约化为
\begin{equation*}
\Gamma^{i}_{j}(r_{\gamma})\Gamma^{i}_{j}(r_{\delta}) = \Gamma^{i}_{j}(r_{\gamma\delta}).
\tag{4.4.11}
\end{equation*}
也就是说，$\Gamma^{i}_{j}$ 是小陪群的一个维数为 $d_{i}\lambda^{i}_{j}$ 的矢量表示。由于 $\sum_{j}d^{2}_{i}\left(\lambda^{i}_{j}\right)^{2} = d^{2}_{i}|\overline{\mathbf{K}}^{i}| > |\overline{\mathbf{K}}^{i}|$，这些表示不可能全都不可约——这与 $\Gamma^{i}_{j}$ 不可约并不矛盾，因为可约的是把它们当作 $\overline{\mathbf{K}}^{i}$ 的表示。
